\section{Introduction}
\label{sec:intro}

Low-rank adaptation (\lora)~\citep{hu2021lora} is the default way to fine-tune large language models on a budget. Nearly every recipe fixes the rank at $r{=}8$ or $r{=}16$, usually with $\alpha=2r$, and rarely justifies the choice. {\bf Is this default a blind spot?} If the rank a task needs varied by an order of magnitude, practitioners would be leaving accuracy on the table for ``hard'' tasks and wasting memory on easy ones.

There are good reasons to suspect such a blind spot. The original \lora work found that very low ranks suffice for GPT-3 adaptation~\citep{hu2021lora}, and fine-tuning has a low intrinsic dimension~\citep{aghajanyan2020intrinsic}. On the other hand, later work reports that code and math need much higher rank in continued pre-training~\citep{biderman2024lora}, and adaptive methods allocate rank unevenly across layers~\citep{zhang2023adalora}. The resulting folk wisdom is simple: sentiment needs rank $\approx 4$, while code needs rank $\geq 32$.

However, two problems make this folk wisdom hard to check. First, rank comparisons are rarely controlled: data budget, optimizer, and evaluation change between studies. Second, under the standard $\alpha \propto r$ scaling, rank and effective step size are coupled~\citep{kalajdzievski2023rslora}. A comparison across ranks at one learning rate therefore mixes ``how much capacity a task needs'' with ``how large a step the learning rate gives at that rank''. Finally, even if rank needs did differ, there is no cheap way to predict them before training.

We run a controlled study to separate these factors. We hold the data budget (4,000 examples, one epoch), the optimizer, and the \lora parameterization fixed. We then sweep eight ranks ($1$ to $128$) on six tasks that range from binary sentiment to math word problems and Python code, on two model sizes of one family (\qsmall with two seeds, \qlarge with one). We add a learning-rate sweep on three tasks to remove the rank--learning-rate confound. We also test a training-free predictor: the effective rank of the full-weight gradient at the pretrained weights. We pre-register four hypotheses: rank needs vary $\geq 8\times$ across tasks (H1), $r{=}8$ costs 3--8\pp on high-rank tasks (H2), the knee rank shrinks with model size (H3), and the gradient statistic predicts the knee with Spearman $\rho > 0.7$ (H4).

Our results mostly refute these hypotheses. At a fixed learning rate, the best rank does span 1 to 128, but the ordering is reversed: \gsm and code are the \emph{lowest}-rank tasks, and both degrade at $r{=}128$ (\gsm falls from 37.3\% at $r{=}1$ to 30.1\%; code falls from 28.8\% to 22.7\%, below the 27.0\% zero-shot baseline). The only significant $r{=}8$ gap after Holm correction is on 6-way \emotion classification ($+2.6$\pp at 0.5B, $+2.2$\pp at 1.5B), and it vanishes when the learning rate is tuned per rank: at learning rate $5\times10^{-4}$, $r{=}8$ reaches 92.8\% and beats $r{=}32$ and $r{=}128$. On \emotion, the learning rate moves $r{=}8$ accuracy by 15\pp, while rank moves it by at most 3\pp at the best learning rate. The gradient statistic separates tasks cleanly (code $\gg$ \gsm $>$ \boolq $>$ \agnews $>$ \emotion $>$ \sst), yet it \emph{anti}-predicts the knee ($\rho=-0.74$, $p=0.014$).

Our contributions are:
\begin{itemize}[leftmargin=*,itemsep=0pt,topsep=0pt]
    \item We conduct a controlled sweep of about 200 \lora adapters across 6 tasks, 8 ranks, and 2 model sizes. We show that at a fixed learning rate, rank needs run opposite to the common intuition: math and code prefer rank 1--16, and large ranks cause collapses of 7--11\pp.
    \item We quantify the cost of the $r{=}8$ default with paired bootstrap tests and cross-fitting across seeds. The cost is at most 2.6\pp, and it is significant only on \emotion.
    \item We show with a learning-rate sweep that the apparent rank effect is mostly a learning-rate effect under $\alpha=2r$. Once the learning rate is tuned, $r{=}8$ is the best rank or within 0.8\pp of it on all three tasks in the sweep.
    \item We evaluate a gradient-spectrum predictor of rank needs and find a significant \emph{negative} correlation with the empirical knee. It lands within a factor of 2 of the knee for 0 of 12 task--model pairs, compared with 6 of 12 for always predicting $r{=}8$.
\end{itemize}

\Secref{sec:related} reviews related work, \secref{sec:method} describes our setup, \secref{sec:results} presents results, \secref{sec:discussion} interprets them and lists limitations, and \secref{sec:conclusion} concludes.
